\begin{center}
{ \LARGE Constitution du Club d'informatique de l'Université d'Ottawa }

%  This should be the date of the last official change, meaning the day of the
%  vote the change.  TODO: review this policy
Dernière modification~: 18 avril 2026
\end{center}

\section*{Article 1 - Nom}

Les noms officiels du club seront ``Club d'informatique de l'Université d'Ottawa'' en français
et ``University of Ottawa Computer Science Club'' en anglais.
L'abréviation semi-officielle ``uOCSClub'' peut également être utilisée.
Aucun autre nom ne sera utilisé dans la publicité ou la représentation du club.

\section*{Article 2 - Mission du club}

% It should be general enough to not dictate us to perform specific activities,
% but to provide a guideline of intention behind our events and activities.

\begin{enumerate}
    % Hackathons, workshops, and other events
    \item Organiser et promouvoir des défis de programmation sociaux et compétitifs~;
    % Managing our Slack channel, partially related to our social events, notifying
    % community of important developments in IT.
    \item Favoriser une communauté de hackers et de programmeurs à l'Université d'Ottawa~;
    %
    \item Offrir des séances de programmation et des ateliers technologiques~;
    \item Organiser des projets d'administration de systèmes menés par la communauté~; et
    \item Être un forum pour l'enseignement et le partage de connaissances en informatique.
\end{enumerate}

\section*{Article 3 - Adhésion}

% This part of the constitution is important.  We need to define what it means to
% be a member of CS Club.

Un individu devient un \emph{membre} reconnu du club lors de son inscription officielle auprès des membres de l'exécutif du club.

Les responsabilités et le mandat de l'exécutif en ce qui concerne l'adhésion au
club sont~:
\begin{enumerate}
    \item S'assurer que l'adhésion est ouverte, accessible et gratuite pour tous les individus~;
    \item S'assurer qu'un minimum de 75\,\% des membres du club sont des membres du SÉUO; et
    \item Exiger que tous les membres de l'exécutif du club soient soit des membres du SÉUO
\end{enumerate}
% Member rights and clauses on disability Sat 11 Oct 2025 13:58:25 EDT
\subsection*{Droits des membres}
Tous les membres du club jouissent des droits suivants~:
\begin{enumerate}
    \item Le droit de voter lors de toutes les élections, amendements et référendums du club~;
    \item Le droit de se porter candidat à l'exécutif, sous réserve des critères d'admissibilité~;
    \item Le droit de proposer des amendements à la Constitution, conformément à l'article 8~;
    \item Le droit d'accéder aux documents de gouvernance du club, y compris les procès-verbaux des réunions, dans un format accessible sur demande~;
    \item Le droit de participer pleinement aux activités du club, indépendamment de tout handicap, avec des accommodements raisonnables fournis sur demande lorsque nécessaire.
\end{enumerate}

\section*{Article 4 - Exécutif}

% Need:
%  2 people to have signing authority
%  to define our council
%  (if CSSA thing happens) official representative
%
Le club sera gouverné par le comité exécutif~; défini ci-dessous~:
\begin{enumerate}
    \item Exécutifs signataires (max 3):
    \begin{enumerate}
        \item Président(e) (ou co-président(e)s)
        \item Trésaurier(ère)
    \end{enumerate}
    \item Exécutifs 
    \begin{enumerate}
        \item VP Logistique
        \item VP Externe
        \item VP CS Games
        \item VP Marketing
    \end{enumerate}
\end{enumerate}

\textit{Il doit y avoir a tout temps au moins deux exécutifs signataires. Les autres postes exécutifs peuvent être vacants.}
\textit{Dans l'éventualité où un poste de l'exécutif devient vacant, ses responsabilités seront partagées entre les autres membres du comité exécutif jusqu'à ce que le poste soit pourvu.}
\textit{Au moins un des exécutifs signataire devra être le/la représentatif(ve) officiel(le) pour l'AÉI }

\section*{Article 5 - Responsabilités de l'exécutif(ve)}


Les responsabilités minimales de l'exécutif(ve) sont les suivantes~:
\begin{enumerate}
    \item Tous les membres de l'exécutif devront~:
    \begin{enumerate}
      \item Organiser au moins un événement par semestre.
    \end{enumerate}
    \item Le/La président(e) devra~:
    \begin{enumerate}
        \item Superviser les autres membres de l'exécutif dans l'accomplissement de leurs
              responsabilités~;
        \item Présider toutes les réunions~;
        \item Avoir l'autorité de signature pour le club~; et
        \item Rédiger les procès-verbaux de toutes les réunions.
    \end{enumerate}

    \item Le/La VP Trésaurier(ère) devra~:  % Addition
    \begin{enumerate}
      \item Maintenir des registres financiers précis de toutes les transactions du club~;
      \item Gérer le compte bancaire du club et l'allocation budgétaire pour les événements, ateliers et compétitions~;
      \item Traiter les demandes de financement (SÉUO, commanditaires, etc.) et les remboursements conformément aux politiques de l'université~;
    \end{enumerate}


    \item Le/La VP Logistique devra~:  % Addition
    \begin{enumerate}
      \item Coordonner la logistique des événements et compétitions, y compris
          les réservations de salles, l'équipement et le transport~;
      \item Gérer l'inventaire tel que les articles promotionnels du club, les
          bannières et le matériel d'événements~;
      \item S'assurer que tous les événements respectent les politiques de
          sécurité et de réservation de l'université~;
      \item Travailler avec le/la VP Finances pour planifier les budgets et les
          achats liés aux événements.
    \end{enumerate}

    \item Le/La VP Externe devra~:
    \begin{enumerate}
        \item Répondre aux courriels du club~;
        \item Contacter des organisations externes pour obtenir des commandites et du financement~;
        \item Assister aux réunions mensuelles de l'AÉÉI (CSSA).
    \end{enumerate}

    \item Le/La VP CS Games devra~:
    \begin{enumerate}
      \item Organiser une session d'introduction annuelle aux CS Games~;
      \item Assurer la coordination des sessions d'entraînement pour les CS Games~;
      \item Coordonner les essais et la logistique pour la compétition annuelle des CS Games.
      \item Assister à la compétition annuelle des CS Games en tant que chef d'équipe.
    \end{enumerate}

    \item Le/La VP Marketing devra~:
    \begin{enumerate}
        \item Publier sur les réseaux sociaux (Discord, Instagram) pour promouvoir les activités du club~;
        \item Contacter d'autres clubs pour des événements promotionnels et des collaborations~;
        \item Répondre aux messages privés sur les plateformes de réseaux sociaux.
    \end{enumerate}
\end{enumerate}



\section*{Article 6 - Officiel(e)s nommé(e)s}

Les exécutifs(ves) peuvent nomer des officiel(e)s pour les aider avec les
tâches administratives du club. Ces officiel(e)s n'ont pas de droit de vote dans les décisions exécutives.
Ces positions peuvent rester vacantes.

Les responsabilités minimales des Officiel(e)s sont les suivantes

\begin{enumerate}
    \item Les conseillers/conseillères nommés devront~:
    \begin{enumerate}
        \item Fournir des conseils et des orientations au comité exécutif sur les affaires du club.
    \end{enumerate}
    \item Le/La directeur/directrice du design nommé devra~:
    \begin{enumerate}
        \item Concevoir le matériel promotionnel du club~;
        \item Assister le/la VP Marketing dans ses fonctions~;
        \item Diriger la conception de tout le matériel destiné au public pour le club.
    \end{enumerate}
    \item L'ambassadeur/ambassadrice de première année nommé devra~:
    \begin{enumerate}
        \item Représenter les étudiants de première année et communiquer leurs préoccupations ou suggestions au comité exécutif~;
        \item Aider à promouvoir les activités du club auprès des nouveaux étudiants et encourager leur participation~;
        \item Aider à organiser des événements axés sur les étudiants de première année, tels que des ateliers d'introduction ou des programmes de mentorat.
    \end{enumerate}
\end{enumerate}


\section*{Article 7 - Meetings}

Le commité exécutif est reponsable d'organiser une assemblée générale par année scolaire et les réunions exécutives nécéssaires.

\begin{enumerate}
    \item L'assemblée générale
        \begin{enumerate}
            \item Doit être annoncée une semaine avant son déroulement
            \item Doit se dérouler à la fin du trimestre d'été
            \item Est ouverte à tous les membres enregistrés de uOCSClub
            \item Peut être en personne, en ligne ou hybride
            \item Si elle est en personne, elle doit se dérouler sur ou autour du campus d'uOttawa.
            \item Les proces verbaux doivent être publiés
        \end{enumerate}
    \item Réunions exécutives
        \begin{enumerate}
            \item Apres l'élection, les nouveaux(elles) exécutifs(ves) décideront comment elles se dérouleront:
            \item Elles n'ont pas besoin d'être annoncées
            \item Les membres peuvent demander quand et où elles se déroulent
            \item Elles sont ouvertes à tous les membres de uOCSClub
        \end{enumerate}
\end{enumerate}

\section*{Article 8 - Élections}

% The election process is for the executive members.

\begin{enumerate}
    \item Les élections pour l'exécutif auront lieu à la fin du semestre d'hiver.
    \begin{enumerate}
        \item Pas avant la dernière semaine de mars~;
        \item Pas après la dernière semaine d'avril.
    \end{enumerate}
    \item La durée du mandat des membres exécutifs(ves) sera de leur élction jusqu'à la fin des élections l'année suivante.
    \item Les positions exécutives sont ouvertes à tous les membres du club qui sont aussi membres de l'SÉUO.
    \item Tout système de vote en ligne utilisé sera administré par une plateforme neutre ou un officiel d'élection désigné par le comité exécutif~;
    \item Toutes les élections seront menées de manière juste et transparente, garantissant une opportunité égale à tous les membres de participer~;
    \item Les élections seront rendues accessibles à tous les membres, y compris la fourniture d'accommodements raisonnables pour les membres en situation de handicap sur demande~;
    \item Dans la mesure du possible, les élections et le matériel connexe seront rendus disponibles en anglais et en français.
    \item Le procédé électoral sera par auto-nomination pour pour une position dans une platforme de sondage en ligne.
    \item Les membres peuvent seulement se nominerpour une position
    \item La période d'auto-nomination durera 5 jours.
    \item Les nominations non-contestées seront élues à la position.
    \item Les nominations contestées procéderons comme suis:
    \begin{enumerate}
        \item Les nominés(es) auront 48h à partir de la fin de la période de nomination mener une campagne pour leur position.
        \item Après la période de campagne, un sondage en ligne sera disponnible pour 24h pour chaque position contestée.
        \item Pendant l'élection, chaque membre du club aura un (1) vote pour chaque position exécutive.
        \item Le/la candidat(e) avec le plus de votes sera élu(e).
        \item Les égalités seront déterminées par une vote majoritaire (2/3)
            des exécutifs(ves) précédents(es) tout en respectant les normes
            pour les procédés justes et d'anti-discrimination du SÉUO.
        \item Après ces procédés, les résultats des élections seront publiés
    \end{enumerate}
    \item Les candidats(es) peuvent retirer leur nomination à tout temps avant leur élection.
    \item Les membres du commité exécutif peuvent disqualifier des candidats
        avec cause par un vote majoritaire (2/3) tant que la cause respecte les
        normes pour les procédés justes et d'anti-discrimination du SÉUO.
    \item Les positions vacantes apres l'élection peuvent être remplies par nomination du nouveau commité exécutif tout en respectant les normes
            pour les procédés justes et d'anti-discrimination du SÉUO.

\end{enumerate}

\section*{Article 9 - Amendements}

\begin{enumerate}
    \item Les amendements à la constitution doivent obtenir une majorité des deux tiers (2/3)
          des voix des membres présents à la réunion;
    \item Au moins la moitié des membres du commité exécutif doivent être présents.
    \item Un amendement à la constitution doit être publié aux membres, qui
          doivent recevoir une copie dactylographiée de l'amendement ainsi que des
          procès-verbaux dactylographiés de la réunion au cours de laquelle l'amendement
          a été adopté, afin de prouver que l'amendement a été approuvé.
\end{enumerate}

\section*{Article 10 - Destitution}
\begin{enumerate}
    \item Tout membre du club qui commet un acte portant atteinte aux
        intérêts du club et de ses membres peut recevoir un avis de
        destitution~;
    \item L'individu visé par la destitution aura le droit de défendre ses
        actions à n'importe quel réunion du club;
    \item Un vote à la majorité des deux tiers (2/3) des membres présents entraînera
        le retrait de l'individu visé du club et la perte
        de tout privilège associé au club.
\end{enumerate}

\section*{Article 11 - Processus de destitution de l'exécutif}

\begin{enumerate}
    \item Dans l'éventualité où le/la président(e) du Club d'informatique de l'Université d'Ottawa s'engage dans des actions nuisibles aux intérêts du club ou manque à ses responsabilités, les membres généraux et les membres de l'exécutif peuvent initier un processus de destitution. Le processus de destitution du/de la président(e) suivra les étapes décrites ci-dessous~:

    \begin{enumerate}
        \item Avis de destitution~: Tout membre ou membre de l'exécutif peut soumettre un avis écrit de destitution au comité exécutif, détaillant les actions ou comportements spécifiques justifiant la destitution. L'avis de destitution doit être présenté lors d'une réunion du comité exécutif, et une copie de l'avis doit être remise au/à la président(e) visé(e)~;

        \item Droit de défense~: Le/La président(e) visé(e) par la destitution aura le droit de se défendre lors d'une réunion spéciale convoquée à cet effet. Cette réunion devra se tenir dans un délai raisonnable après la soumission de l'avis de destitution. Le/La président(e) pourra présenter sa défense et fournir des preuves ou des témoins en sa faveur~;

        \item Vote de destitution~: Après que le/la président(e) a présenté sa défense, un vote de destitution sera tenu. Une majorité des deux tiers (2/3) des voix des membres présents à la réunion, incluant les membres généraux et les membres de l'exécutif, sera requise pour destituer le/la président(e) de son poste~;

        \item Perte de privilèges~: Si le vote de destitution est adopté, le/la président(e) visé(e) sera immédiatement retiré(e) de son poste et perdra tous les privilèges associés à la présidence du club. Les fonctions et responsabilités du/de la président(e) seront transférées au/à la vice-président(e) ou à un autre membre de l'exécutif désigné, jusqu'à l'élection d'un(e) nouveau/nouvelle président(e).
    \end{enumerate}


    \item Si un membre de l'exécutif, autre que le/la président(e), s'engage dans des actions nuisibles aux intérêts du club ou manque à ses responsabilités, les membres généraux et les autres membres de l'exécutif peuvent initier un processus de destitution. Le processus de destitution des membres de l'exécutif suivra une procédure similaire à celle décrite pour la destitution du/de la président(e), avec les modifications nécessaires pour tenir compte du poste spécifique en question.
\end{enumerate}

\section*{Article 12 - Clause d'agence}

% This would need to be updated to include CSSA or other overseeing organizations.
\begin{enumerate}
    \item Le Club d'informatique de l'Université d'Ottawa n'est pas un mandataire de
          l'Université d'Ottawa et ses opinions et actions ne représentent pas celles
          de l'Université d'Ottawa.

      \item Le Club d'informatique de l'Université d'Ottawa n’est pas un agent
          du Syndicat étudiant de l’Université d’Ottawa (SÉUO) et ses opinions
          ne sont pas représentatives de celles du SÉUO, sauf indication
          contraire du SÉUO.

      \item Le Club d'informatique de l'Université d'Ottawa n’est pas un agent
          de l'Association des Étudiants en informatique de l’Université
          d’Ottawa (AÉI) et ses opinions ne sont pas représentatives de celles
          du AÉI, sauf indication contraire du AÉI.

\end{enumerate}

% Include discussion on club meetings

\section*{Article 13 - Équité et accessibilité}
% Added this to stay compliant with UOSU  Sat 11 Oct 2025 13:55:52 EDT
\begin{enumerate}
    \item Le club affirme son engagement envers l'équité, la diversité et l'inclusion dans toutes ses opérations, activités et événements~;
    \item Le club ne discriminera pas sur la base de la race, du genre, de la sexualité, du handicap, de la religion, de la langue ou de tout autre motif protégé par les lois de l'Ontario~;
    \item Toutes les réunions, élections et activités du club seront rendues accessibles aux membres en situation de handicap, avec des accommodements raisonnables fournis sur demande~;
    \item Le club s'efforcera de fournir des communications et des services bilingues en anglais et en français, reconnaissant la nature bilingue de l'Université d'Ottawa.
\end{enumerate}
